% Generated draft: requires mathematical and source review.
\documentclass[11pt]{article}
\usepackage[margin=25mm]{geometry}
\usepackage{fontspec}
\setmainfont{Latin Modern Roman}
\usepackage{amsmath,amssymb,mathtools,unicode-math}
\setmathfont{Latin Modern Math}
\usepackage{enumitem,needspace,xurl,hyperref}
\setlength{\parindent}{0pt}
\setlength{\parskip}{5pt}
\setlength{\emergencystretch}{3em}
\newcommand{\N}{\mathbb{N}}
\newcommand{\Z}{\mathbb{Z}}
\newcommand{\Q}{\mathbb{Q}}
\newcommand{\R}{\mathbb{R}}
\newcommand{\C}{\mathbb{C}}
\newcommand{\F}{\mathbb{F}}
\newcommand{\A}{\mathbb{A}}
\newcommand{\PP}{\mathbb{P}}
\newcommand{\E}{\mathbb{E}}
\newcommand{\eps}{\varepsilon}
\newcommand{\dd}{\,\mathrm{d}}
\begin{document}
\section*{UnsolvedMath ID: 800850}
\textbf{Source problem code:} catalogue entry
\subsection{Definitions and mathematical statement}
For a simple undirected graph \(H\) on \(n\) vertices, define its Lov\'asz number by
\[
\vartheta(H)=\max\left\{\sum_{i,j=1}^{n}X_{ij}: X\in\R^{n\times n},\ X=X^{\mathsf T},\ X\succeq0,\ \operatorname{tr}(X)=1,\ X_{ij}=0\text{ whenever }\{i,j\}\in E(H)\right\}.
\]
Here \(X\succeq0\) means that \(X\) is positive semidefinite. Under the usual simple undirected convention, a circulant graph has vertex set \(\Z/n\Z\) and a set \(S\subseteq(\Z/n\Z)\setminus\{0\}\) satisfying \(S=-S\), with distinct vertices \(x,y\) adjacent exactly when \(y-x\in S\).

Let \(G_n\) denote a graph on \(n\) vertices drawn from the source's intended random dense circulant graph model. The supplied conjecture is
\[
\E\,\vartheta(G_n)=(1+o(1))\sqrt n,
\]
or equivalently
\[
\lim_{n\to\infty}\frac{\E\,\vartheta(G_n)}{\sqrt n}=1,
\]
where the expectation is over the intended graph distribution and the limit is along its admissible vertex counts.

The supplied record does not define that distribution or its admissible vertex counts. In particular, it does not specify the law of \(S\), the edge density parameter, independence assumptions, or any parity restrictions. The circulant definition above states the usual simple undirected convention; its applicability to the intended source model requires confirmation. Consequently, the displayed asymptotic is a conditional transcription of the conjectured target, rather than a fully specified probabilistic statement.

The source attributes the conjecture to numerical observations and comparison with \(G(n,1/2)\), the random simple graph in which each unordered pair of distinct vertices is independently an edge with probability \(1/2\). This comparison is supplied motivation and does not itself specify the circulant model.
\subsection{Short English statement}
The conjecture predicts that the expected Lov\'asz number of a random dense circulant graph on \(n\) vertices is asymptotic to \(\sqrt n\). It is motivated by supplied numerical observations and a comparison with \(G(n,1/2)\). The random circulant distribution is missing from the record, so the precise probabilistic formulation remains unresolved.
\subsection{Sources}
\begin{itemize}
\item Local UnsolvedMath registry: permanent ID \(800850\), source problem code \texttt{catalogue entry}; title: \emph{Randomstrasse101 Open Problems (2025) --- Conjecture 2}. The supplied registry attributes the problem to Bandeira et al., records proposed year 2026, and labels its status as open. These are supplied registry data, not independently verified assessments.
\item Bandeira et al., \emph{Randomstrasse101: Open Problems of 2025} (2026), Conjecture 2. Supplied source URL: \url{https://arxiv.org/abs/2603.29571}. The background records access on 2026-10-16 and extraction by \texttt{arxiv\_latex}; these provenance details are supplied and were not independently verified here.
\item Supplied upstream revision: \texttt{d371ee917ac9647fbb06cd6fc529deec6f3dda82}. The record was created and updated on 2026-10-16. No separate literature-verification assessment is supplied.
\end{itemize}
\end{document}
